\section*{Hoofdstuk \ref{ch:b3:lebesgue} --- De
Lebesgue-integraal}

\begin{solution}{exo:b3:lebesgue:1}
(a) Is $f$ niet-dalend, dan is $\{f > t\}$ gelijk aan
$\varnothing$, $\R$ of een halfrechte $\intoo a{+\infty}$ of
$\intco a{+\infty}$: in elk geval borel; en voor niet-stijgende
$f$ net zo. Een afgeleide: $f'(x) = \lim_n n\bigl(f(x + \frac1n)
- f(x)\bigr)$ is een puntsgewijze limiet van \omterm{def:b3:topology:continuity}{continue} (en dus
\omterm{def:b3:lebesgue:measurable}{meetbare}) functies: \cref{prop:b3:lebesgue:stability}(c).

(b) $\{f > t\} = \bigcup_{q \in \Q,\, q > t}\{f > q\}$: zijn de
rationale niveaus \omterm{def:b3:lebesgue:measurable}{meetbaar}, dan alle niveaus, en de halfrechten
brengen $\mathcal B(\R)$ voort.
\end{solution}

\begin{solution}{exo:b3:lebesgue:2}
Eerste: $(1 + x/n)^n \nearrow \eu^x$ voor $x \geq 0$, dus gaat de
integrand puntsgewijs naar $\eu^{-x}\cos x$; en voor $n \geq 2$ is
$(1 + x/n)^n \geq (1 + x/2)^2$, wat de \omterm{def:b3:lebesgue:l1}{integreerbare} dominant $(1
+ x/2)^{-2}$ geeft. Gedomineerde convergentie:
\[
\lim_n\int_0^\infty\frac{\cos x}{(1 + x/n)^n}\dd x
= \int_0^\infty \eu^{-x}\cos x\,\dd x
= \operatorname{Re}\int_0^\infty\eu^{-(1 - \iu)x}\dd x
= \operatorname{Re}\frac{1}{1 - \iu} = \frac12 .
\]
Tweede: substitueer $u = x^n$ (een $\mathcal C^1$-bijectie van
$\intoo01$):
\[
\int_0^1\frac{nx^{n-1}}{1 + x}\dd x = \int_0^1\frac{\dd u}{1 +
u^{1/n}} \longrightarrow \int_0^1\frac{\dd u}{2} = \frac12,
\]
via de gedomineerde convergentie: voor $u \in \intoo01$ is
$u^{1/n} \to 1$, en de integrand is begrensd door $1$ op een
ruimte van eindige \omterm{def:b3:measure:measure}{maat}.
\end{solution}

\begin{solution}{exo:b3:lebesgue:3}
(a) $f_n = n\,\mathbf 1_{\intoo0{1/n}}$: puntsgewijs is $\liminf
f_n = 0$, en $\int f_n = 1$: dus $0 < 1$.

(b) Hoogte: $n\mathbf 1_{\intoo0{1/n}}$; breedte: $\frac1n\mathbf
1_{\intoo0n}$; verschuiving: $\mathbf 1_{\intoo n{n+1}}$. Alle
gaan puntsgewijs naar $0$ met $\int = 1$. In elk geval faalt de
hypothese van de \emph{dominatie}: $\sup_nf_n$ is respectievelijk
$\approx 1/x$ nabij $0$, ongeveer een \omterm{def:b3:lebesgue:l1}{niet-integreerbaar} constant
profiel, en van het type $\mathbf 1_{\intoo1\infty}$ --- nooit
\omterm{def:b3:lebesgue:l1}{integreerbaar}.

(c) ``$\int\limsup f_n \geq \limsup\int f_n$'' faalt voor de
verschuivende bult: linkerlid $0$, rechterlid $1$. En
``$\limsup\int \leq \int\limsup$'' is dezelfde uitspraak. Ook
Fatou voor de $\limsup$ met omgekeerde $\leq$ (``omgekeerde
Fatou'') vergt een dominant --- dezelfde bult is het
tegenvoorbeeld.
\end{solution}

\begin{solution}{exo:b3:lebesgue:4}
(a) Voor $x > 0$ is $\frac1{\eu^x - 1} = \frac{\eu^{-x}}{1 -
\eu^{-x}} = \sum_{n\geq1}\eu^{-nx}$, dus $\frac{x}{\eu^x - 1} =
\sum_{n\geq1}x\eu^{-nx}$, een reeks van niet-negatieve \omterm{def:b3:lebesgue:measurable}{meetbare
functies}: \cref{cor:b3:lebesgue:additivity} staat term voor term
integreren toe:
\[
\int_0^\infty\frac{x\,\dd x}{\eu^x - 1}
= \sum_{n\geq1}\int_0^\infty x\eu^{-nx}\dd x
= \sum_{n\geq1}\frac1{n^2} = \frac{\pi^2}6
\]
($\int_0^\infty x\eu^{-nx}\dd x = n^{-2}$ met partiële
integratie; het Bazelprobleem uit het volume van bachelorjaar 2,
of \cref{exo:b3:hilbert:5} verderop).

(b) Op $\intoo01$ is $-x\ln x \geq 0$, dus is $x^{-x} = \eu^{-x\ln
x} = \sum_k\frac{(-x\ln x)^k}{k!}$ een reeks met niet-negatieve
termen: opnieuw verwisselen. Substitutie $x = \eu^{-u/(k+1)}$
geeft
\[
\int_0^1(-x\ln x)^k\dd x
= \int_0^\infty\Bigl(\frac{u}{k+1}\Bigr)^{k}
\eu^{-\frac{ku}{k+1}}\;\frac{\eu^{-\frac u{k+1}}}{k+1}\,\dd u
= \frac{1}{(k+1)^{k+1}}\int_0^\infty u^k\eu^{-u}\dd u
= \frac{k!}{(k+1)^{k+1}} .
\]
Bijgevolg is $\int_0^1x^{-x}\dd x = \sum_{k\geq0}\frac{1}{(k+1)
^{k+1}} = \sum_{n\geq1}n^{-n}$: de droom van de tweedejaars,
streng bewezen.
\end{solution}

\begin{solution}{exo:b3:lebesgue:5}
(a) Markov: $a\,\mathbf 1_{\{f \geq a\}} \leq f$; integreren geeft
$\mu(\{f \geq a\}) \leq \frac1a\int f$. Is $\int f = 0$, dan is
$\mu(\{f \geq \frac1n\}) = 0$ voor elke $n$, en is $\{f > 0\} =
\bigcup_n\{f \geq \frac1n\}$ een nulverzameling.

(b) $\mu(\{\abs f = \infty\}) \leq \mu(\{\abs f \geq n\}) \leq
\frac1n\int\abs f \to 0$.

(c) Neem $A = \{f > 0\}$: dan is $\int f^+\dd\mu = \int_Af\,\dd\mu
= 0$, dus $f^+ = 0$ \omterm{def:b3:lebesgue:l1}{bijna overal} volgens (a); en net zo $f^- = 0$
\omterm{def:b3:lebesgue:l1}{bijna overal}.
\end{solution}

\begin{solution}{exo:b3:lebesgue:6}
(a) $\mathbf 1_\Q$: overal discontinu, dus niet
\omterm{def:b3:lebesgue:l1}{riemann-integreerbaar} (\cref{thm:b3:lebesgue:riemann}); haar
lebesgue-integraal is $\lambda(\Q) = 0$. Thomae: \omterm{def:b3:topology:continuity}{continu} in elk
irrationaal punt (bij gegeven $\varepsilon$ hebben slechts eindig
veel rationale getallen in $\intcc01$ een noemer $\leq
1/\varepsilon$; mijd ze met een kleine \omterm{def:b3:topology:topology}{omgeving}) en discontinu in
de rationale punten (\omterm{ex:b3:lebesgue:gamma}{dichtheid} van de irrationale getallen): dus
\omterm{def:b3:lebesgue:l1}{bijna overal} \omterm{def:b3:topology:continuity}{continu}, \omterm{def:b3:lebesgue:l1}{riemann-integreerbaar}, met integraal $0$ (ze
verdwijnt \omterm{def:b3:lebesgue:l1}{bijna overal}). $\mathbf 1_K$ met $K$ een dikke
cantorverzameling: de discontinuïteitsverzameling is $\partial K =
K$ (gesloten met leeg inwendige), van \omterm{def:b3:measure:measure}{maat} $\frac12 > 0$: dus
\emph{niet} \omterm{def:b3:lebesgue:l1}{riemann-integreerbaar} --- en toch
\omterm{def:b3:lebesgue:l1}{lebesgue-integreerbaar} met integraal $\lambda(K) = \frac12$.

(b) $\int_{k\pi}^{(k+1)\pi}\frac{\abs{\sin x}}x\dd x \geq
\frac1{(k+1)\pi}\int_{k\pi}^{(k+1)\pi}\abs{\sin x}\dd x =
\frac{2}{(k+1)\pi}$: de reeks divergeert. Convergentie van de
oneigenlijke integraal: voor $A > \pi$ is
\[
\int_\pi^A\frac{\sin x}x\dd x = \Bigl[\frac{-\cos
x}x\Bigr]_\pi^A - \int_\pi^A\frac{\cos x}{x^2}\dd x,
\]
en beide termen convergeren als $A \to \infty$ ($\frac1{x^2}$ is
\omterm{def:b3:lebesgue:l1}{integreerbaar}): voorwaardelijke convergentie zonder absolute
\omterm{def:b3:lebesgue:l1}{integreerbaarheid}.
\end{solution}

\begin{solution}{exo:b3:lebesgue:7}
Dominatie: $\abs{\partial_t(\eu^{-x^2}\cos(tx))} =
\abs{x\eu^{-x^2}\sin(tx)} \leq x\eu^{-x^2}$, \omterm{def:b3:lebesgue:l1}{integreerbaar} en
onafhankelijk van $t$: \cref{thm:b3:lebesgue:paramdiff} is dus
globaal van toepassing, en
\[
F'(t) = -\int_0^\infty x\eu^{-x^2}\sin(tx)\,\dd x
= \Bigl[\tfrac12\eu^{-x^2}\sin(tx)\Bigr]_0^\infty
- \frac t2\int_0^\infty\eu^{-x^2}\cos(tx)\dd x
= -\frac t2F(t)
\]
(partiële integratie met $\dd v = x\eu^{-x^2}\dd x$). De lineaire
differentiaalvergelijking geeft $F(t) = F(0)\eu^{-t^2/4}$; en met
$F(0) = \frac{\sqrt\pi}2$ (\cref{pb:b3:lebesgue:1}) reproduceert
de gaussische functie zichzelf onder deze cosinustransformatie.
\end{solution}

\begin{solution}{exo:b3:lebesgue:8}
De integraal convergeert: nabij $0$ gaat de integrand naar $b - a$
(begrensd), en in het oneindige neemt ze af als $\eu^{-ax}$. Leg
$a$ vast en bekijk $I(b) = \int_0^\infty\frac{\eu^{-ax} -
\eu^{-bx}}x\dd x$ als functie van $b \in \intco a{+\infty}$. Voor
alle $b \geq a$ is $\abs{\partial_b(\text{integrand})} =
\eu^{-bx} \leq \eu^{-ax}$, en $\int_0^\infty\eu^{-ax}\dd x =
\frac1a < \infty$: een \omterm{def:b3:lebesgue:l1}{integreerbare} dominant. Dus geeft
\cref{thm:b3:lebesgue:paramdiff} dat $I'(b) =
\int_0^\infty\eu^{-bx}\dd x = \frac1b$, met $I(a) = 0$:
\[
I(b) = \int_a^b\frac{\dd t}t = \ln\frac ba .
\]
(Gelijkwaardig, langs de route van de hint: de integrand is
$\int_a^b\eu^{-xt}\dd t \geq 0$ en de verwisseling is het \omterm{def:b3:topology:continuity}{continue}
analogon van \cref{cor:b3:lebesgue:additivity}, dat wil zeggen
Tonelli --- bewezen in \cref{ch:b3:product}; de parameterroute
blijft binnen dit hoofdstuk.)
\end{solution}

\begin{solution}{exo:b3:lebesgue:9}
$\nu(\varnothing) = 0$; en voor disjuncte $(A_n)$ is $f\mathbf
1_{\bigsqcup A_n} = \sum_nf\mathbf 1_{A_n}$ (puntsgewijs, alle
termen $\geq 0$), zodat \cref{cor:b3:lebesgue:additivity} de
$\sigma$-additiviteit geeft. De formule $\int g\,\dd\nu = \int
gf\,\dd\mu$: voor $g = \mathbf 1_A$ is ze de definitie van $\nu$;
voor elementaire $g$ volgt ze uit de lineariteit; en voor \omterm{def:b3:lebesgue:measurable}{meetbare}
$g \geq 0$ neem je elementaire $s_n \nearrow g$
(\cref{thm:b3:lebesgue:approximation}): dan is $s_nf \nearrow gf$,
en gaat de monotone convergentie aan beide kanten tot de limiet
over. (Die roltrap ``indicator $\to$ elementair $\to$ monotone
convergentie'' is de standaardmachine van de theorie.)
\end{solution}

\begin{solution}{exo:b3:lebesgue:10}
(a) Uit $\abs{\sin(tx)} \leq \abs tx$ volgt
$\abs{\frac{\sin(tx)}{x(1+x^2)}} \leq \frac{\abs t}{1+x^2}$: de
integraal convergeert, en op $\abs t \leq T$ levert de dominant
$\frac{T}{1+x^2}$ de \omterm{def:b3:topology:continuity}{continuïteit}
(\cref{thm:b3:lebesgue:paramcont}). Differentiëren:
$\abs{\partial_t} = \abs{\frac{\cos(tx)}{1+x^2}} \leq
\frac1{1+x^2}$, \omterm{def:b3:lebesgue:l1}{integreerbaar}, dus $f'(t) =
\int_0^\infty\frac{\cos(tx)}{1+x^2}\dd x$ voor alle $t$. Een
tweede differentiatie zou het integreren van
$\frac{x\sin(tx)}{1 + x^2}$ vergen, waarvan de absolute waarde
zich in het oneindige gedraagt als $\frac{\abs{\sin(tx)}}x$:
\emph{niet} \omterm{def:b3:lebesgue:l1}{integreerbaar} --- er is geen dominant en
\cref{thm:b3:lebesgue:paramdiff} kan geen tweede keer worden
toegepast.

(b) Nemen we $f'(t) = \frac\pi2\eu^{-t}$ voor $t > 0$ aan, dan is
$f'$ wegens de oneven aard van $f$ even, dus $f'(t) =
\frac\pi2\eu^{-\abs t}$ --- en dat is \emph{niet differentieerbaar
in $0$}: $f$ is $\mathcal C^1$ maar niet $\mathcal C^2$, wat
bevestigt dat de geblokkeerde tweede differentiatie geen technisch
toeval was. Voor $t > 0$ is de oneigenlijke integraal
$\int_0^\infty\frac{x\sin(tx)}{1+x^2}\dd x$ gelijk aan $-f''(t) =
\frac\pi2\eu^{-t}$ (differentieer de aangenomen formule waar dat
geoorloofd is, dus op $\intoo0\infty$) --- een oneigenlijke,
niet-lebesguewaarde.
\end{solution}

\begin{solution}{exo:b3:lebesgue:11}
(a) Zij $g_n = (f - f_n)^+$: dan is $0 \leq g_n \leq f$ (want
$f_n \geq 0$), $g_n \to 0$ \omterm{def:b3:lebesgue:l1}{bijna overal}, en $f$ een \omterm{def:b3:lebesgue:l1}{integreerbare}
dominant: dus $\int g_n \to 0$ (gedomineerde convergentie). En
omdat $\abs{f_n - f} = 2(f - f_n)^+ - (f - f_n)$, is
\[
\int\abs{f_n - f} = 2\int g_n - \Bigl(\int f - \int
f_n\Bigr) \longrightarrow 0 + 0 .
\]

(b) De schuivende bult $f_n = \mathbf 1_{\intcc n{n+1}}$ heeft
$f_n \to 0$ \omterm{def:b3:lebesgue:l1}{bijna overal} en $\int f_n = 1 \not\to 0$: zonder de
convergentie van de integralen faalt de $L^1$-convergentie (en de
hypothese ook). Het voorbeeld met wisselend teken: $f_n \to 0$ in
elke $x \neq 0$ en $\int f_n = 0$, maar $\int\abs{f_n} = 2$: voor
rijen met wisselend teken gaat de stelling werkelijk over controle
van het type $\abs{f_n}$, en de positiviteit werd precies gebruikt
in $g_n \leq f$.

(c) Dichtheden voldoen aan $\int p_n = 1 = \int p$: de hypothese
van (a) geldt automatisch, dus dwingt $p_n \to p$ \omterm{def:b3:lebesgue:l1}{bijna overal} af
dat $\norm{p_n - p}_{L^1} \to 0$ --- en daarmee de convergentie
van de kansen $\int_Ap_n \to \int_Ap$, \emph{uniform} over alle
\omterm{def:b3:lebesgue:measurable}{meetbare} $A$ ($\abs{\int_A(p_n - p)} \leq \norm{p_n - p}_1$):
Scheffé maakt van puntsgewijze convergentie van dichtheden
convergentie van de wetten in totale variatie.
\end{solution}

\begin{solution}{exo:b3:lebesgue:12}
(a) Op $\intoo0n$ stijgt $\varphi_n(x) = n\log(1 - \frac xn)$ in
$n$ naar $-x$ (de afbeelding $t \mapsto \frac{\log(1 - xt)}{t}$
daalt als $t = \frac1n \downarrow 0$; of ontwikkel: $\varphi_{n}'
\geq 0$ in $n$ via $\log(1-u) + \frac{u}{1-u} \geq 0$). Dus $(1 -
\frac xn)^n\eu^{x/2}\mathbf 1_{x<n} \nearrow \eu^{-x/2}$, en de
monotone convergentie geeft
\[
\lim_n\int_0^n\Bigl(1 - \frac xn\Bigr)^n\eu^{x/2}\dd x =
\int_0^\infty\eu^{-x/2}\dd x = 2 .
\]

(b) Substitueer $u = x^n$ (dus $x = u^{1/n}$ en $n x^{n-1}\dd x =
\dd u$):
\[
\int_0^1\frac{nx^{n-1}}{1 + x}\dd x =
\int_0^1\frac{\dd u}{1 + u^{1/n}} .
\]
Voor $u \in \intoo01$ is $u^{1/n} \to 1$, dus gaat de integrand
naar $\frac12$, gedomineerd door $1$: de limiet is $\frac12$
(gedomineerde convergentie). (De massa van $nx^{n-1}$ concentreert
zich bij $x = 1$, waar $\frac1{1+x} = \frac12$: de substitutie
maakt die concentratie zichtbaar.)

(c) Puntsgewijs is $(1 + x/n)^n \nearrow \eu^x$ en $x^{1/n} \to 1$
(voor $x > 0$): de integrand gaat naar $\eu^{-x}$. Dominant voor
$n \geq 2$: op $\intoc01$ is $x^{-1/n} \leq x^{-1/2}$ en $(1 +
x/n)^{-n} \leq 1$, wat de \omterm{def:b3:lebesgue:l1}{integreerbare} grens $x^{-1/2}$ geeft; en
op $\intoo1\infty$ is $x^{-1/n} \leq 1$ en $(1 + x/n)^n \geq 1 +
\binom n2\frac{x^2}{n^2} \geq 1 + \frac{x^2}4$, wat de
\omterm{def:b3:lebesgue:l1}{integreerbare} grens $\frac{4}{4 + x^2}$ geeft. Gedomineerde
convergentie:
\[
\lim_n\int_0^\infty\frac{\dd x}{(1 + x/n)^nx^{1/n}} =
\int_0^\infty\eu^{-x}\dd x = 1 .
\]
\end{solution}

\begin{solution}{pb:b3:lebesgue:1}
\textbf{1.} $A$ is $\mathcal C^1$ wegens de hoofdstelling van de
integraalrekening en de kettingregel: $A'(t) =
2\eu^{-t^2}\int_0^t\eu^{-x^2}\dd x$. Voor $B$ is
$\partial_t\bigl(\frac{\eu^{-t^2(1+x^2)}}{1+x^2}\bigr) =
-2t\,\eu^{-t^2(1+x^2)}$, \omterm{def:b3:topology:continuity}{continu} en begrensd op $[t_0, T] \times
\intcc01$ voor elke $0 < t_0 < T$ (een begrensde dominant op een
ruimte van eindige \omterm{def:b3:measure:measure}{maat} volstaat): $B$ is dus $\mathcal C^1$ op
$\intoo0{+\infty}$ met
\[
B'(t) = -2t\int_0^1 \eu^{-t^2(1+x^2)}\dd x
= -2\eu^{-t^2}\int_0^1 t\,\eu^{-(tx)^2}\dd x
= -2\eu^{-t^2}\int_0^t\eu^{-u^2}\dd u = -A'(t)
\]
(substitutie $u = tx$).

\textbf{2.} $A(0) = 0$ en $B(0) = \int_0^1\frac{\dd x}{1+x^2} =
\frac\pi4$. $A + B$ heeft op $\intoo0{+\infty}$ afgeleide nul en
is \omterm{def:b3:topology:continuity}{continu} in $0$ ($B$ via de dominant $\frac1{1+x^2}$ en
\cref{thm:b3:lebesgue:paramcont}): dus $A + B \equiv \frac\pi4$.
Als $t\to+\infty$: $0 \leq B(t) \leq \eu^{-t^2} \to 0$, en $A(t)
\to \bigl(\int_0^\infty\eu^{-x^2}\dd x\bigr)^2$ (monotone
convergentie van de binnenste integraal).

\textbf{3.} Bijgevolg is $\bigl(\int_0^\infty\eu^{-x^2}\dd
x\bigr)^2 = \frac\pi4$, dus $\int_0^\infty\eu^{-x^2}\dd x =
\frac{\sqrt\pi}2$, en wegens de pariteit $G = \sqrt\pi$. Verder is
$\Gamma(\frac12) = \int_0^\infty x^{-1/2}\eu^{-x}\dd x
\overset{x = u^2}{=} 2\int_0^\infty\eu^{-u^2}\dd u = \sqrt\pi$.

\textbf{4.} Voor $t > 0$ is $\abs{\eu^{-tx}\frac{\sin x}x} \leq
\eu^{-tx}$, \omterm{def:b3:lebesgue:l1}{integreerbaar}. Voor $t = 0$ de oneigenlijke
convergentie: op $[\pi, A]$ is
\[
\int_\pi^A\frac{\sin x}x\dd x
= \Bigl[-\frac{\cos x}x\Bigr]_\pi^A -
\int_\pi^A\frac{\cos x}{x^2}\dd x,
\]
en beide termen convergeren als $A \to \infty$; en nabij $0$ zet
de integrand zich \omterm{def:b3:topology:continuity}{continu} voort met de waarde $1$.

\textbf{5.} Op $[t_0, +\infty)$ met $t_0 > 0$ is
$\abs{\partial_t\bigl(\eu^{-tx}\tfrac{\sin x}x\bigr)} =
\eu^{-tx}\abs{\sin x} \leq \eu^{-t_0x}$, \omterm{def:b3:lebesgue:l1}{integreerbaar}: dus is
\cref{thm:b3:lebesgue:paramdiff} op elk zulk interval van
toepassing, en daarmee op heel $\intoo0{+\infty}$:
\[
F'(t) = -\int_0^\infty\eu^{-tx}\sin x\,\dd x
= -\operatorname{Im}\int_0^\infty\eu^{-(t - \iu)x}\dd x
= -\operatorname{Im}\frac{1}{t - \iu} = -\frac{1}{1 + t^2}.
\]

\textbf{6.} $\abs{F(t)} \leq \int_0^\infty\eu^{-tx}\dd x =
\frac1t \to 0$. Integreren van $F' = -\frac1{1+t^2}$ geeft $F(t) =
C - \arctan t$, en $t \to \infty$ dwingt $C = \frac\pi2$ af: dus
$F(t) = \frac\pi2 - \arctan t$ op $\intoo0{+\infty}$.

\textbf{7.} Integreer partieel op $[A, R]$ met $\sin x = (-\cos
x)'$ en laat $R \to \infty$ gaan:
\[
\int_A^{\infty}\eu^{-tx}\frac{\sin x}x\dd x
= \frac{\cos A\;\eu^{-tA}}{A}
- \int_A^\infty \cos x\,\Bigl(\frac tx +
\frac1{x^2}\Bigr)\eu^{-tx}\dd x .
\]
Met $\abs{\cos} \leq 1$: de eerste term is $\leq \frac1A$; en de
integraal is hoogstens $\int_A^\infty t\eu^{-tx}\frac{\dd x}x +
\int_A^\infty\frac{\dd x}{x^2} \leq \frac1A\int_0^\infty
t\eu^{-tx}\dd x + \frac1A = \frac2A$. Totaal: $\leq \frac 3A$,
uniform voor $t \in [0, 1]$ (het geval $t = 0$ inbegrepen). Nu is
\[
\abs{F(t) - D} \leq
\Bigl|\int_0^A(\eu^{-tx} - 1)\,\frac{\sin x}x\,\dd x\Bigr| +
\frac6A .
\]
Op $[0, A]$ is $\abs{\eu^{-tx} - 1} \leq tx \leq tA$ en
$\abs{\frac{\sin x}x} \leq 1$, dus is de eerste term hoogstens
$tA^2$. Kies $A$ met $\frac6A < \varepsilon$, en daarna $t <
\varepsilon/A^2$: dan is $\abs{F(t) - D} < 2\varepsilon$.

\textbf{8.} Bijgevolg is $D = \lim_{t\to0^+}F(t) =
\lim_{t\to0^+}\bigl(\frac\pi2 - \arctan t\bigr) = \frac\pi2$.

\textbf{9.} Partiële integratie ($u = \sin^2x$, $v' = x^{-2}$):
\[
\int_0^\infty\frac{\sin^2x}{x^2}\dd x
= \Bigl[-\frac{\sin^2x}{x}\Bigr]_0^\infty +
\int_0^\infty\frac{2\sin x\cos x}{x}\dd x
= \int_0^\infty\frac{\sin 2x}{x}\dd x = D = \frac\pi2
\]
(substitueer $u = 2x$ in de laatste stap; de randtermen
verdwijnen: $\sin^2 x/x \to 0$ aan beide uiteinden).

\textbf{10.} Partiële integratie ($u = 1 - \cos x$, $v' =
x^{-2}$): $\int_0^\infty\frac{1 - \cos x}{x^2}\dd x =
\int_0^\infty\frac{\sin x}x\dd x = \frac\pi2$. Consistentie: $1 -
\cos x = 2\sin^2\frac x2$, en de substitutie $x = 2u$ maakt van
$\int\frac{2\sin^2(x/2)}{x^2}\dd x$ de integraal
$\int\frac{\sin^2u}{u^2}\dd u$: de twee berekeningen stroken.

\textbf{11.} $\int_0^\infty\frac{\sin(ax)}x\dd x = \frac\pi2$ voor
elke $a > 0$ (substitueer $u = ax$: de integraal is
schaalinvariant); en $\int_\R\eu^{-ax^2}\dd x = \sqrt{\pi/a}$
(substitueer $u = \sqrt a\,x$). Schalen in het argument laat de
integraal van Dirichlet onveranderd en deelt de gaussische
integraal door $\sqrt a$.

\textbf{12.} Een dominant die voor alle $t \in [0,1]$ geldt, moet
$\sup_{t\in[0,1]}\abs{\eu^{-tx}\frac{\sin x}x} = \abs{\frac{\sin
x}x}$ domineren, en dat is niet \omterm{def:b3:lebesgue:l1}{integreerbaar}
(\cref{exo:b3:lebesgue:6}): de gedomineerde convergentie kan $t =
0$ dus niet oversteken. Het alternatief van vraag 7 --- staarten
uniform klein in de parameter, en het \omterm{def:b3:topology:compact}{compacte} deel met de
gedomineerde convergentie --- is het standaardpatroon van de
``uniforme \omterm{def:b3:lebesgue:l1}{integreerbaarheid}'', en keert terug zodra
voorwaardelijke convergentie een limietverwisseling ontmoet.

\textbf{13.} Is $g = \log f$ convex, dan is $f = \exp\circ g$
convex (exp is convex en stijgend: $f(\lambda x + (1-\lambda)y)
\leq \eu^{\lambda g(x) + (1-\lambda)g(y)} \leq \lambda f(x) +
(1-\lambda)f(y)$, met de laatste stap wegens de convexiteit van
exp tussen de punten $g(x)$ en $g(y)$). Producten en affiene
substituties: logaritmen maken er sommen en affiene substituties
van convexe functies van. Hölder (met $\frac1p + \frac1q = 1$):
voor $\int\abs u^p = \int\abs v^q = 1$ geeft Young dat $\abs{uv}
\leq \frac{\abs u^p}p + \frac{\abs v^q}q$; integreren geeft
$\int\abs{uv} \leq 1$, en het algemene geval volgt uit de
homogeniteit. Pas het dan voor $\lambda \in \intoo01$ met $p =
\frac1\lambda$ toe op de ontbinding
\[
t^{\lambda x + (1-\lambda)y - 1}\eu^{-t}
= \bigl(t^{x-1}\eu^{-t}\bigr)^{\lambda}
\bigl(t^{y-1}\eu^{-t}\bigr)^{1-\lambda} :
\qquad
\Gamma(\lambda x + (1{-}\lambda)y) \leq
\Gamma(x)^\lambda\,\Gamma(y)^{1-\lambda} .
\]

\textbf{14.} Voor een convexe $g$ stijgt de helling van een koorde
met haar eindpunten (de driekoordenongelijkheid). Vergelijking van
de koorden over $[n-1, n]$, $[n, n+x]$ en $[n, n+1]$ geeft
\[
\log(n-1) = \frac{g(n) - g(n-1)}1 \leq
\frac{g(n+x) - g(n)}x \leq \frac{g(n+1) - g(n)}1 = \log n,
\]
en vermenigvuldigen met $x > 0$ geeft de bewering.

\textbf{15.} Met $f(n) = (n-1)!$ (recursie vanaf $f(1) = 1$) en
$f(n + x) = x(x+1)\cdots(x+n-1)\,f(x)$ luidt vraag 14:
\[
(n-1)^x\,(n-1)! \;\leq\; x(x+1)\cdots(x+n-1)\,f(x)
\;\leq\; n^x\,(n-1)! .
\]
De bovengrens herschrijft zich als $f(x) \leq
\frac{n!\,n^x}{x(x+1)\cdots(x+n)}\cdot\frac{x+n}n$, en de
ondergrens op rang $n+1$ als $f(x) \geq
\frac{n!\,n^x}{x(x+1)\cdots(x+n)}$. De correctiefactor
$\frac{x+n}n \to 1$: de insluiting dwingt
\[
f(x) = \lim_n\frac{n!\,n^x}{x(x+1)\cdots(x+n)}
\qquad (x \in \intoc01)
\]
af, een uitdrukking die niet van $f$ afhangt: eenduidigheid op
$\intoc01$, en dus overal wegens de recursie. En omdat $\Gamma$
aan alle drie de hypothesen voldoet (vraag 13), is $f = \Gamma$ en
geldt de formule van Gauss --- voor alle $x > 0$, want beide leden
voldoen aan dezelfde recursie.

\textbf{16.} Nabij $0$ is de integrand $\sim t^{x-1}$,
\omterm{def:b3:lebesgue:l1}{integreerbaar} precies wanneer $x > 0$; en nabij $1$ symmetrisch
met $y$. Partiële integratie op $\intcc\varepsilon{1-\varepsilon}$
en $\varepsilon \to 0$ (de randtermen verdwijnen voor $x, y > 0$):
\[
B(x{+}1, y) = \Bigl[-t^x\frac{(1-t)^y}y\Bigr]_0^1 +
\frac xy\int_0^1t^{x-1}(1-t)^y\,\dd t
= \frac xy\bigl(B(x, y) - B(x{+}1, y)\bigr),
\]
met $(1-t)^y = (1-t)^{y-1}(1 - t)$; oplossen geeft $B(x+1, y) =
\frac{x}{x+y}B(x, y)$. En $B(1, y) = \int_0^1(1-t)^{y-1}\dd t =
\frac1y$.

\textbf{17.} $f(1) = \frac1y\cdot\frac{\Gamma(1+y)}{\Gamma(y)} =
1$; $f(x+1) = \frac{x}{x+y}B(x,y)\cdot\frac{(x+y)\Gamma(x+y)}
{\Gamma(y)} = x\,f(x)$; en $f$ is logaritmisch convex in $x$ als
product van de logaritmisch convexe $B(\cdot, y)$ (Hölder op de
ontbinding $t^{(\lambda x_1 + (1-\lambda)x_2)-1}(1-t)^{y-1} =
(\cdots)^\lambda(\cdots)^{1-\lambda}$, als in vraag 13) en
$\Gamma(\cdot + y)$ (affiene verschuiving). Bohr--Mollerup: $f =
\Gamma$, dat wil zeggen $B(x, y) =
\frac{\Gamma(x)\Gamma(y)}{\Gamma(x+y)}$.

\textbf{18.} Met $t = \sin^2\theta$ is $\dd t =
2\sin\theta\cos\theta\,\dd\theta$ en $t^{-1/2}(1-t)^{-1/2} =
\frac1{\sin\theta\cos\theta}$:
\[
B\Bigl(\frac12, \frac12\Bigr) =
\int_0^{\pi/2}2\,\dd\theta = \pi
= \frac{\Gamma(\frac12)^2}{\Gamma(1)}
\quad\Longrightarrow\quad
\Gamma\Bigl(\frac12\Bigr) = \sqrt\pi .
\]
Deel I bereikte dezelfde constante door naar een parameter te
differentiëren en twee functies naar hun limieten te laten
wedijveren; hier maakte de convexiteit alleen het probleem zo
stijf dat er nog maar één waarde overbleef. Analyse door beweging
tegenover analyse door vorm.

\textbf{19.} $g(1) = \frac{2^0}{\sqrt\pi}\Gamma(\frac12)\Gamma(1)
= 1$. Recursie:
\[
g(x+1) = \frac{2^{x}}{\sqrt\pi}\,
\Gamma\Bigl(\frac{x+1}2\Bigr)\Gamma\Bigl(\frac x2 + 1\Bigr)
= \frac{2^{x}}{\sqrt\pi}\cdot\frac x2\,
\Gamma\Bigl(\frac x2\Bigr)\Gamma\Bigl(\frac{x+1}2\Bigr)
= x\,g(x) .
\]
Logaritmische convexiteit: product van $\eu^{(x-1)\log2}$
(logaritmisch affien) en twee affiene herparametriseringen van de
logaritmisch convexe $\Gamma$. Bohr--Mollerup geeft $g = \Gamma$;
en met $x = 2z$: $\Gamma(2z) =
\frac{2^{2z-1}}{\sqrt\pi}\Gamma(z)\Gamma(z + \frac12)$.

\textbf{20.} Uit $\Gamma(\frac12) = \sqrt\pi$ en de recursie volgt
$\Gamma(n + \frac12) = (n - \frac12)(n - \frac32)\cdots\frac12\,
\sqrt\pi = \frac{(2n-1)(2n-3)\cdots1}{2^n}\sqrt\pi =
\frac{(2n)!}{4^nn!}\sqrt\pi$ (vul het oneven product met de even
getallen aan). Asymptotiek: vraag 14 met $g = \log\Gamma$ en $x =
\frac12$ geeft $\sqrt{n-1} \leq
\frac{\Gamma(n+\frac12)}{\Gamma(n)} \leq \sqrt n$, zodat de
verhouding tot $\sqrt n$ tussen $\sqrt{1 - \frac1n}$ en $1$ wordt
ingeklemd.

\textbf{21.} Uit vraag 20 volgt $(2n)! =
\frac{4^nn!}{\sqrt\pi}\Gamma(n + \tfrac12)$, dus
\[
\binom{2n}n = \frac{(2n)!}{(n!)^2}
= \frac{4^n\,\Gamma(n+\frac12)}{\sqrt\pi\;n!}
= \frac{4^n}{\sqrt\pi}\cdot
\frac{\Gamma(n+\frac12)}{n\,\Gamma(n)}
\sim \frac{4^n}{\sqrt\pi}\cdot\frac{\sqrt n}{n}
= \frac{4^n}{\sqrt{\pi n}} .
\]
Numeriek is $4^{10}/\sqrt{10\pi} = 1048576/5.6050 \approx
187078.6$, tegenover $\binom{20}{10} = 184756$: verhouding
$0.9876$ --- de fout is $O(1/n)$, bij $n = 10$ nog zichtbaar.

\textbf{22.} Gelijkwaardigheden: $\Gamma(\frac12) = \sqrt\pi
\Leftrightarrow G = \sqrt\pi$ (de substitutie $x = u^2$ uit vraag
3) $\Leftrightarrow B(\frac12, \frac12) = \pi$ (de formule van
Euler); en de verdubbeling in $z = n$ \emph{is} de gesloten vorm
van $\Gamma(n + \frac12)$, die op het hellingslemma na \emph{is}
de schatting van de centrale binomiaalcoëfficiënt. De
parametertechniek (Delen I--II) berekent limieten van bewegende
grootheden en is onmisbaar zodra er een echte vervorming in het
spel is (de integraal van Dirichlet heeft geen convexiteitsbewijs);
de convexiteitstechniek berekent niets maar verbiedt alles --- ze
blinkt uit in eenduidigheid en functionaalvergelijkingen (Gauss,
Euler, Legendre in drie streken), waar het differentiëren in
rekenwerk zou verzuipen. Een volleerd analyticus draagt beide.

\textbf{23.} Met $t = \sin^2\theta$ is $\dd t =
2\sin\theta\cos\theta\,\dd\theta = 2\,t^{1/2}(1-t)^{1/2}\,
\dd\theta$, dus
\[
W_p = \int_0^1 t^{p/2}\,
\frac{\dd t}{2\,t^{1/2}(1-t)^{1/2}}
= \frac12\int_0^1 t^{\frac{p+1}2 - 1}(1-t)^{\frac12 - 1}\dd t
= \frac12\,B\Bigl(\frac{p+1}2, \frac12\Bigr).
\]
De bètarecursie met $x = \frac{n+1}2$ en $y = \frac12$ geeft
\[
W_{n+2} = \frac12\,
\frac{(n+1)/2}{(n+2)/2}\,B\Bigl(\frac{n+1}2, \frac12\Bigr)
= \frac{n+1}{n+2}\,W_n .
\]
Vanaf $W_0 = \frac\pi2$ en $W_1 = 1$:
\[
W_{2n} = \frac\pi2\prod_{k=1}^n\frac{2k-1}{2k},
\qquad
W_{2n+1} = \prod_{k=1}^n\frac{2k}{2k+1} .
\]
Omdat $\sin^{n+1} \leq \sin^n$ op $\intcc0{\pi/2}$, is de rij
$(W_n)$ niet-stijgend, zodat
\[
1 \geq \frac{W_{2n+1}}{W_{2n}} \geq
\frac{W_{2n+1}}{W_{2n-1}} = \frac{2n}{2n+1}
\longrightarrow 1 .
\]
Maar de gesloten vormen geven
\[
\frac{W_{2n+1}}{W_{2n}} = \frac2\pi
\prod_{k=1}^n\frac{(2k)^2}{(2k-1)(2k+1)}
= \frac2\pi\prod_{k=1}^n\frac{4k^2}{4k^2-1},
\]
en $n \to \infty$ levert het product van Wallis. (Via de formule
van Euler is $W_p = \frac{\sqrt\pi}2\,
\Gamma(\frac{p+1}2)/\Gamma(\frac p2 + 1)$: Wallis is de gaussische
integraal in weer een ander kostuum.)

\textbf{24.} De afgeleide van de integrand naar $b$ is
$-2x\,\eu^{-x^2}\sin(2bx)$, gedomineerd door $2\abs
x\,\eu^{-x^2} \in L^1(\R)$, uniform in $b$: dus is $\Phi$
$\mathcal C^1$ met
\[
\Phi'(b) = -\int_{-\infty}^{+\infty}
2x\,\eu^{-x^2}\sin(2bx)\,\dd x .
\]
Partiële integratie met $u = \sin(2bx)$ en $\dd v =
-2x\,\eu^{-x^2}\dd x$ (dus $v = \eu^{-x^2}$) laat de randtermen
verdwijnen en geeft
\[
\Phi'(b) = -2b\int_{-\infty}^{+\infty}
\eu^{-x^2}\cos(2bx)\,\dd x = -2b\,\Phi(b).
\]
Bijgevolg is $\bigl(\Phi(b)\,\eu^{b^2}\bigr)' = 0$ en $\Phi(b) =
\Phi(0)\,\eu^{-b^2} = \sqrt\pi\,\eu^{-b^2}$ volgens Deel I. Op de
normering na zegt dat dat de fouriergetransformeerde van
$\eu^{-x^2}$ weer een gaussische functie is --- het vaste punt
waarop de inversietheorie van \cref{ch:b3:fouriertransform}
draait.

\textbf{25.} Voor $0 < a < b$ en $x > 0$ is
\[
0 \leq \frac{\eu^{-ax} - \eu^{-bx}}{x}
= \int_a^b \eu^{-sx}\,\dd s \leq (b - a)\,\eu^{-ax},
\]
dus convergeert de integraal $I(a)$ (in de zin van Lebesgue); en
de puntsgewijze grens laat ook zien dat de integrand zich in $x =
0$ \omterm{def:b3:topology:continuity}{continu} voortzet met de waarde $b - a$. Leg $b$ vast; op
$\intco{a_0}{+\infty}$ is de afgeleide van de integrand naar $a$
gelijk aan $-\eu^{-ax}$, gedomineerd door $\eu^{-a_0x} \in
L^1(\intoo0{+\infty})$, dus is $I$ $\mathcal C^1$ op $\intoo0b$
met
\[
I'(a) = -\int_0^{+\infty}\eu^{-ax}\dd x = -\frac1a,
\qquad\text{en dus}\qquad
I(a) = -\log a + c .
\]
De tweezijdige grens geeft $0 \leq I(a) \leq (b-a)/a \to 0$ als $a
\to b^-$, dus $c = \log b$ en $I(a) = \log\frac ba$. De ophefbare
singulariteit is \emph{in $0$} nodig: elke term $\eu^{-ax}/x$
heeft afzonderlijk een divergente (logaritmische) integraal nabij
$0$, en alleen de opheffing van de eerste orde $\eu^{-ax} -
\eu^{-bx} = O(x)$ maakt het verschil daar \omterm{def:b3:lebesgue:l1}{integreerbaar}; in het
oneindige is elke term op zichzelf al onschadelijk.
\end{solution}
